\section{Data}\label{sec:data}

\subsection{Policy documents and PCSI}

The policy data come from the corpus assembled for the companion paper \citep{SharmaPCSI2026}: 518 IP and patent policy documents from 150 U.S.\ research universities, dated 1944--2025, which form 480 institution-year document records and 87,159 sentences. Each sentence is classified as restrictive, neutral, or supportive in how it presents rights, obligations, authority, and assistance, and PCSI is a document's mean sentence-level score, so higher values mean more supportive communication toward inventors. The companion paper describes how the measure is built and validated. For the supporting annual panel, each observed document is treated as the policy in force until the next one appears.

\subsection{Technology-transfer outcomes}

Technology-transfer outcomes come from the AUTM Licensing Activity Survey for 1991--2023. The raw data contain 253 reporting identifiers, which we harmonize to 149 institutions. This step matters. Over the period, 101 institutions report under more than one identifier, and treating each identifier as a separate institution would split institution fixed effects and break lagged variables whenever the identifier changes. Every analysis is therefore keyed to the harmonized institution.

We study five principal outcomes. Four are counts, entered as $\log(1+y)$: licenses and options executed, invention disclosures received, new U.S.\ patent applications filed, and U.S.\ patents issued. The fifth is the filing margin,
\begin{equation}
F_{it}=\log(1+\text{Applications}_{it})-\log(1+\text{Disclosures}_{it}),
\label{eq:filingmargin}
\end{equation}
an institution-year log-difference version of a filing ratio \citep{ChoudhryPonzio2020}. The AUTM survey does not link individual disclosures to individual applications, so $F_{it}$ is not an invention-level conversion probability. Each of the five outcomes is estimated separately in the stacked event study below.

Two series needed extra checking. In the FY2023 survey, the total licenses-and-options field no longer equals licenses issued plus options issued, as it does in almost every record from 2005 to 2022, and its 2023 median is far below the sum of the two parts. We therefore use licenses issued plus options issued for 2023, which keeps the definition of the licensing series consistent (Appendix Table~\ref{tab:licfield}). Startups formed are reported only from 1994, and institutions self-report them under a definition tied to licensing university technology. Short coverage, a changing share of zeros, and institutions entering and leaving the survey make this series harder to compare over time than the five principal outcomes, so we use startups only as supplementary evidence (Appendix~\ref{app:startups}). We leave out licensing income because its reported units change across survey years.

Table~\ref{tab:descriptives} summarizes the linked panel, which has 3,507 institution-year observations. PCSI is available for 2,564 of them, but most of those annual values are carried forward from an earlier document. That is one reason the event design looks at documented changes in the sequence of policies instead of treating every carried-forward annual value as a new observation.

\begin{table}[htbp]\centering
\caption{Linked AUTM--PCSI panel, 1991--2023}\label{tab:descriptives}
\begin{threeparttable}\small\begin{tabular}{lrrr}\toprule
Variable & $N$ & Mean & SD \\\midrule
PCSI (policy in force) & 2,564 & 0.440 & 0.063 \\
Invention disclosures & 3,507 & 122.4 & 157.7 \\
New patent applications & 3,507 & 70.6 & 113.7 \\
U.S. patents issued & 3,351 & 33.0 & 50.3 \\
Licenses and options executed$^{a}$ & 3,489 & 40.0 & 69.2 \\
Research expenditure (\$ millions) & 3,507 & 333.0 & 478.1 \\
Licensing FTEs & 3,401 & 5.5 & 7.3 \\
Startups formed$^{b}$ & 3,245 & 4.6 & 7.1 \\
\midrule
Institution-years & \multicolumn{3}{r}{3,507} \\
Harmonized institutions (AUTM reporting identifiers) & \multicolumn{3}{r}{149 (253)} \\
Share of annual PCSI values carried forward & \multicolumn{3}{r}{0.893} \\
Within-institution share of PCSI variance & \multicolumn{3}{r}{0.30} \\
\bottomrule\end{tabular}
\begin{tablenotes}[flushleft]\footnotesize
\item \textit{Notes:} Institution-year observations. PCSI is the policy-in-force score from the companion measurement paper, carried forward from the most recent observed document. $^{a}$For FY2023 the survey's total field no longer equals licenses plus options issued, so 2023 values use licenses plus options issued. $^{b}$Supplementary outcome, reported from 1994 (Appendix~\ref{app:startups}). Licensing income is not used because reported units are inconsistent across survey years.
\end{tablenotes}\end{threeparttable}\end{table}

\subsection{Revisions and the revision sample}

We identify revisions from consecutive \emph{observed} policy documents in the full corpus. A threshold revision occurs when PCSI changes between two consecutive observed documents by
\[
|\Delta\mathrm{PCSI}|>0.03.
\]
A revision is \emph{upward} when $\Delta\mathrm{PCSI}>0.03$, a move toward more supportive communication, and \emph{downward} when $\Delta\mathrm{PCSI}<-0.03$, a move toward more restrictive communication. The threshold is about half the cross-document standard deviation of PCSI. Appendix Table~\ref{tab:thresholduniverse} shows how many revisions nearby thresholds would produce.

We build the set of revisions from the full sequence of observed documents, not from the annual AUTM panel, because a policy document can appear in a year with no AUTM observation. For this set, the year of the successor document serves as the mechanical revision date. We use it to find threshold changes and to screen comparison institutions. The corpus is not a complete legal history of every version a university ever adopted, so we do not assume that the year a document was observed is the year it was adopted.

For that reason we keep two questions apart: whether an event is mechanically eligible, and whether its documents are usable. At the baseline threshold we observe 126 revisions at 78 AUTM-linked institutions. The mechanical screen drops a revision when the same institution has another threshold revision within the event window $[E-4,E+5]$, or when the institution has no year with any core AUTM outcome before the event year $E$, or none from $E$ through $E+5$. The screen looks only at whether outcome data exist, never at their values or at estimated effects. Sixty-two revisions pass. Twenty-five of these had already been reviewed by hand in an earlier version of the paper, and we then reviewed the predecessor and successor documents for all 37 remaining pairs.

Documentary review grades each pair on two dimensions. For \emph{comparability}, a pair is C when both documents are the same governing IP or patent policy, P when they are only partly comparable (an excerpt and a full policy, say, or a campus policy and a governing-board policy), and N when they cannot be compared. For \emph{timing}, a pair is A when an adoption or revision date is stated and no intervening version is documented, or when the two documents are at most two years apart; B when an adoption or effective date is stated but the earlier revision history is incomplete; and C when intervening versions exist or no usable date can be established. For events we keep, the event year is the documented adoption or revision year whenever one is available and differs from the year the successor document was observed. The full 126-revision set used for the mechanical screen and for excluding controls keeps the dates of the observed document sequence.

Of the 37 newly reviewed pairs, nine meet both requirements (C or P comparability and A or B timing), and 28 fail at least one. Adding those nine to the earlier events gives the preferred sample: 34 revisions at 32 institutions, eight upward and 26 downward. The main specification keeps an event if an outcome is observed in the revision year, even when the AUTM series ends right after it. Late events therefore have little post-revision data, so we also report a stricter sample that requires at least two years with core outcomes for the revising institution in event times 0 through $+5$. That guarantees at least one observed year after the revision. Applied to every event, the rule leaves 29 revisions, seven upward and 22 downward.

\begin{table}[htbp]\centering
\caption{Construction of the revision sample}\label{tab:revsample}
\begin{threeparttable}\small
\begin{tabular}{p{11.3cm}r}\toprule
Step & Count \\\midrule
Observed policy-document records (institution-years) in the companion corpus & 480 \\
Revisions with $|\Delta\mathrm{PCSI}|>0.03$ at AUTM-linked institutions & 126 \\
Pass mechanical non-overlap and outcome-support screen & 62 \\
\quad reviewed by hand in an earlier version of the paper & 25 \\
\quad additional mechanically eligible pairs reviewed in the second pass & 37 \\
Newly usable C/P pairs with A/B timing from the second pass & 9 \\
Preferred fully documentary-reviewed sample & 34 \\
\quad upward / downward & 8 / 26 \\
\quad distinct revising institutions & 32 \\
Stricter post-support sensitivity sample & 29 \\
\quad upward / downward & 7 / 22 \\
\bottomrule
\end{tabular}
\begin{tablenotes}[flushleft]\footnotesize
\item \textit{Notes:} Mechanical eligibility is determined without using estimated outcome effects. Documentary usability requires C or P comparability and timing tier A or B. The 64 revisions failing the mechanical screen cannot enter the preferred stacked event study regardless of documentary classification. The stricter sample requires at least two treated-institution core-outcome panel years in event times 0--5.
\end{tablenotes}
\end{threeparttable}
\end{table}
